%; whizzy document
% latex beamer presentation.
% platex, latex-beamer でコンパイルすることを想定。 

%     Tokyo Debian Meeting resources
%     Copyright (C) 2006-2007 Junichi Uekawa , Nobuhiro Iwamatsu

%     This program is free software; you can redistribute it and/or modify
%     it under the terms of the GNU General Public License as published by
%     the Free Software Foundation; either version 2 of the License, or
%     (at your option) any later version.

%     This program is distributed in the hope that it will be useful,
%     but WITHOUT ANY WARRANTY; without even the implied warranty of
%     MERCHANTABILITY or FITNESS FOR A PARTICULAR PURPOSE.  See the
%     GNU General Public License for more details.

%     You should have received a copy of the GNU General Public License
%     along with this program; if not, write to the Free Software
%     Foundation, Inc., 51 Franklin St, Fifth Floor, Boston, MA  02110-1301 USA

% 実行順番
% sudo  ~/bin/usb-macbook-ir.c &
% real presentation (shell-command (concat "DISPLAY=:0.1 xpdf -fullscreen " (replace-regexp-in-string "tex$" "pdf"(buffer-file-name)) "&"))
% DISPLAY=:0.1 xpdf -fullscreen 

\documentclass[cjk,dvipdfmx]{beamer}
%\usetheme{Warsaw}
\usetheme{Boadilla}

%  preview (shell-command (concat "xpdf " (replace-regexp-in-string "tex$" "pdf"(buffer-file-name)) "&"))
%  presentation (shell-command (concat "xpdf -fullscreen " (replace-regexp-in-string "tex$" "pdf"(buffer-file-name)) "&"))

%http://www.naney.org/diki/dk/hyperref.html
%日本語EUC系環境の時
\AtBeginDvi{\special{pdf:tounicode EUC-UCS2}}
%シフトJIS系環境の時
%\AtBeginDvi{\special{pdf:tounicode 90ms-RKSJ-UCS2}}

\title{東京エリア Debian 勉強会}
\subtitle{勉強会紹介}
\author{岩松 信洋 iwamatsu@debian.or.jp\\IRC nick: iwamatsu}
\date{2007年03月17日}
\logo{\includegraphics[width=8cm]{image200607/openlogo-light.eps}}

% commandline環境を定義。画面入出力についてはcommandline環境
% で表記する
\newenvironment{commandline}%
{\VerbatimEnvironment
  \begin{Sbox}\begin{minipage}{15cm}\begin{fontsize}{7.3}{7.3} \begin{BVerbatim}}%
{\end{BVerbatim}\end{fontsize}\end{minipage}\end{Sbox}
  \setlength{\fboxsep}{8pt}\fbox{\TheSbox}}

% 三択問題用
\newcounter{santakucounter}
\newcommand{\santaku}[5]{%
\addtocounter{santakucounter}{1}
\frame{\frametitle{問題\arabic{santakucounter}. #1}
%問題\arabic{santakucounter}. #1
\begin{minipage}[t]{0.7\hsize}
 \begin{itemize}
 \item  A #2\\
 \item  B #3\\
 \item  C #4\\
 \end{itemize}
\end{minipage}
}
\frame{\frametitle{問題\arabic{santakucounter}. #1}
%問題\arabic{santakucounter}. #1
\begin{minipage}[t]{0.7\hsize}
\begin{itemize}
\item  A #2\\
\item  B #3\\
\item  C #4\\
\end{itemize}
\end{minipage}
\begin{minipage}[t]{0.2\hsize}
答えは:


\vspace{1cm}

{\huge \hspace{1cm}#5}
\end{minipage}}
}

\begin{document}
\frame{\titlepage{}}

\section{東京エリア Debian 勉強会}

\begin{frame}
 \frametitle{東京エリア Debian 勉強会とは？}
\end{frame}

\begin{frame}
 \frametitle{東京エリア Debian 勉強会とは？}
\begin{itemize}
  \item<1-> Debian Official Developerである上川 純一 さんが発起人.
  \item<2-> 2005年 1月から開始.
  \item<3-> 月に1回コンスタントに勉強会を行っている.
  \item<4-> 場所は荻窪にある公民館（あんさんぶる荻窪）.
\end{itemize}
\end{frame}

\section{勉強会 の目的}

\begin{frame}
 \frametitle{勉強会 の目的}
\begin{itemize}
  \item<1-> 現状は ML と IRC と WEB だけが情報源. 情報交換に限りがある.
  \item<2-> 実際にface to face で話し合える場所がない.
  \item<3-> 情報が断片的で、まとまったドキュメントがない.
\end{itemize}
\end{frame}

\section{勉強会 の目的}
\begin{frame}
 \frametitle{勉強会 の目的}
\begin{itemize}
  \item<1-> 定期的に集まれる場所を作ろう.
  \item<2-> 勉強会を行うときは必ず資料を作成する.
  \item<3-> もちろん、資料は DFSG Free . Webサイトで公開 .
\footnote{http://tokyodebian.alioth.debian.org}
  \item<4-> 最終的にはみんなガリガリパッケージを作れるように.
\end{itemize}
\end{frame}

\section{勉強会 の内容}
\begin{frame}
 \frametitle{勉強会 の内容}
\begin{itemize}
  \item<1-> Debian Weekly News Quiz
  \item<2-> 毎回 Debian に関する マニアックなお題
    \begin{itemize}
	\item Debian Policy について
	\item dpatch / quilt / dbs の使いかた
	\item Debconf 現地からの報告
	\item などなど
    \end{itemize}
  \item<3-> もちろん Debian に関するユーザー向けの話題もやってます
    \begin{itemize}
	\item Macbook on Debian
	\item Debian sid へいざない
	\item BTS のしかた
	\item などなど
    \end{itemize}
  \item<4-> GPG キーサインパーティ
  \item<5-> 宴会という名の交流会
\end{itemize}
\end{frame}

\section{勉強会 を行ってからの結果}
\begin{frame}
 \frametitle{勉強会 を行ってからの結果}
\begin{itemize}
  \item<1-> Debian Package メンテナーが増えた.
    \begin{itemize}
	\item 小林さん
	\item やまねさん
	\item 三塚さん
	\item 岩松
    \end{itemize}	
  \item<2-> Debian Developer 予備軍が増えた.
  \item<3-> Debian に関する資料が増えた.
  \item<4-> Debian での日本語の環境がよくなりつつある.
  \item<5-> Debian JP Project の正式勉強会になった.
  \item<6-> 出張 Debian 勉強会を行うようになった.
    \begin{itemize}
	\item 北海道
	\item 大阪
    \end{itemize}	
\end{itemize}
\end{frame}


\section{今後の予定しているお題}
\begin{frame}
 \frametitle{今後の予定しているお題}
\begin{itemize}
  \item Debian インストールから実戦投入まで
  \item Debconf 7 報告会
  \item xxxxxxxx
\end{itemize}
\end{frame}

\section{今後の開催予定}
\begin{frame}
 \frametitle{今後の開催予定}
\begin{itemize}
  \item 2007年 4月 xx日

	18:00 から 荻窪にて開催予定
  \item 2007年 5月 xx日

	18:00 から 荻窪にて開催予定
\end{itemize}
\end{frame}

\begin{frame}
 \frametitle{Web サイト}
\begin{itemize}
  \item 東京エリア Debian 勉強会 Web サイト

	\url{http://tokyodebian.alioth.debian.org}
\end{itemize}
\end{frame}

\end{document}
